% Part: intuitionistic-logic
% Chapter: propositions-as-types
% Section: sequent-natural-deduction

\documentclass[../../../include/open-logic-section]{subfiles}

\begin{document}

\olfileid{int}{pty}{snd}

\olsection{相继式自然演绎}

我们用 $\Gamma \Sequent !A$ 表示存在一个以 $\Gamma$ 为未解除假设、以 $!A$ 为结论的自然演绎推导；若 $\Gamma$ 为空，则写作 $\Sequent !A$。

我们用 $\Gamma, !A_1, \dots, !A_n$ 表示 $\Gamma \cup \{!A_1, \dots,
!A_n\}$，并用 $\Gamma, \Delta$ 表示 $\Gamma \cup \Delta$。

注意，如果我们有 $\Gamma \Sequent !A \land !A$，即存在一个以 $\Gamma$ 为未解除假设、$!A \land !A$ 为终结公式的推导，那么在底部应用 $\Elim{\land}$，就能得到一个具有相同未解除假设且以 $!A$ 为结论的推导。换言之，如果 $\Gamma \Sequent !A \land !B$，那么 $\Gamma \Sequent !A$。

\begin{prooftree}
  \Axiom$\Gamma \fCenter !A \land !B$
  \RightLabel{$\Elim{\land}$}
  \UnaryInf$\Gamma \fCenter !A$
  \DisplayProof\qquad\bottomAlignProof
  \Axiom$\Gamma \fCenter !A \land !B$
  \RightLabel{$\Elim{\land}$}
  \UnaryInf$\Gamma \fCenter !B$
\end{prooftree}

标签 $\Elim{\land}$ 提示了其与自然演绎中同名规则的关系。

同样地，假设我们有 $\Gamma, !A \Sequent !B$，这意味着我们有一个以 $\Gamma, !A$ 为未解除假设、以~$!B$ 为终结公式的推导。如果我们应用 $\Intro{\lif}$ 规则，就能得到一个以 $\Gamma$ 为未解除假设、$!A \lif !B$ 为终结公式的推导，即 $\Gamma \Sequent !A \lif !B$。请注意，这使得假设的解除变得更为显式。

\begin{prooftree}
  \Axiom $\Gamma, !A \fCenter !B$
  \RightLabel{$\Intro{\lif}$}
  \UnaryInf $\Gamma \fCenter !A \lif !B$
\end{prooftree}

我们可以用同样的方式从其他规则得出推论，具体阐述如下：
\begin{gather*}
  \Axiom $\Gamma \fCenter !A$
  \Axiom $\Delta \fCenter !B$
  \RightLabel{$\Intro{\land}$}
  \BinaryInf $\Gamma,\Delta \fCenter !A \land !B$
  \DisplayProof\\
  \Axiom $\Gamma \fCenter !A \land !B$
  \RightLabel{$\Elim{\land}_1$}
  \UnaryInf $\Gamma \fCenter !A$
  \DisplayProof
  \qquad
  \Axiom $\Gamma \fCenter !A \land !B$
  \RightLabel{$\Elim{\land}_2$}
  \UnaryInf $\Gamma \fCenter !B$
  \DisplayProof
  \\
  \Axiom $\Gamma \fCenter !A$
  \RightLabel{$\Intro{\lor}_1$}
  \UnaryInf $\Gamma \fCenter !A \lor !B$
  \DisplayProof\qquad
  \Axiom $\Gamma \fCenter !B$
  \RightLabel{$\Intro{\lor}_2$}
  \UnaryInf $\Gamma \fCenter !A \lor !B$
  \DisplayProof\\
  \Axiom $\Gamma \fCenter !A \lor !B$
  \Axiom $\Delta, !A \fCenter !C$
  \Axiom $\Delta', !B \fCenter !C$
  \RightLabel{$\Elim{\lor}$}
  \TrinaryInf $\Gamma, \Delta, \Delta' \fCenter !C$
  \DisplayProof
  \\
  \Axiom $\Gamma, !A \fCenter !B$
  \RightLabel{$\Intro{\lif}$}
  \UnaryInf $\Gamma \fCenter !A \lif !B$
  \DisplayProof
  \qquad
  \Axiom $\Delta \fCenter !A \lif !B$
  \Axiom $\Gamma \fCenter !A$
  \RightLabel{$\Elim{\lif}$}
  \BinaryInf $\Gamma, \Delta \fCenter !B$
  \DisplayProof
  \\
  \Axiom $\Gamma \fCenter \lfalse$
  \RightLabel{$\FalseInt$}
  \UnaryInf $\Gamma \fCenter !A$
  \DisplayProof
\end{gather*}

任何假设本身就是一个从~$!A$ 到~$!A$ 的推导，也就是说，我们总是有 $!A \Sequent !A$。

\begin{prooftree}
  \AxiomC{$ $}
  \UnaryInfC{$!A \fCenter !A$}
\end{prooftree}

综合来看，这些规则可以被视为一个关于存在何种自然演绎推导的演算。它们也可以被视为自然演绎的一种记法变体，在这种变体中，每一步不仅记录了所推导出的公式，也记录了推导它所依赖的未解除的假设。

\begin{prooftree}
  \Axiom$!A \fCenter !A$
  \UnaryInf $!A \fCenter !A \lor (!A \lif \lfalse)$
  \Axiom $!B \fCenter !B$
  \BinaryInf$!A, !B\lif \fCenter \lfalse$
  \UnaryInf$(!B \fCenter !A \lif \lfalse$
  \UnaryInf$( !B \fCenter !A \lor (!A \lif \lfalse)$
  \Axiom$(!B \fCenter !B$
  \BinaryInf$(!B \fCenter \lfalse$
  \UnaryInf$\fCenter !B \lif \lfalse$
\end{prooftree}

其中 $!B$ 是 $(!A \lor (!A \lif \lfalse))\lif \lfalse$ 的简写。

\end{document}
